% page 604 of the facsimile (image p-603 of the scan)
% block 144.8 x 233.3 mm ; left margin 8.0 mm
\pgc{-12.700mm}{0.000mm}{
\l{\cel{3}596.}
\l{}
\l{\sou{triumviro}. (en\cel{1}\sou{la Roma antiqua}). Stat-oficiero adminis-}
\l{\cel{3}trala, alta-ranga, qua, lor la republiko, funcionis}
\l{\cel{3}kun du kolegi. - (\sou{anke metaf.}) - DEFIRS.}
\l{}
\l{\sou{triviala}. Sen originaleso, infra-sorta e pasable gro-}
\l{\cel{3}siera. - DEFIRS.}
\l{}
\l{\sou{tro.}\cel{3}Plu (multe) kam konvenas, kam necesa, kam bezo-}
\l{\cel{3}nata, kam demandita. - FI.}
\l{}
\l{\sou{trofeo}. I. (\sou{epoki antiqua}). Asemblajo ek la spoliuro,}
\l{\cel{3}la raptajo de enemiko vinkita, quin onu suspendis a}
\l{\cel{3}la trunko di arboro, kom monumento pri sua vinko.-}
\l{\cel{3}II. - (\sou{Nun-epoke})\cel{2}Grupo de armi, de standardi,}
\l{\cel{3}asalte konquestita de la enemiko, quan onu suspen-}
\l{\cel{3}das kom memorigivo pri sua vinko. - DEFIRS.}
\l{}
\l{\sou{trofoplasmo}. (\sou{histologio}). Parto di la hialoplasmo,o}
\l{\cel{3}citoplasmo, amorfa, qua pleas rolo pure nutrala -}
\l{\cel{3}nomo di la substanco fibretoza fundamentala di la}
\l{\cel{3}neurono. - DEFIRS.}
\l{}
\l{\sou{trogo}. I. Vazo ligna quan uzas la masonisti por dilutar}
\l{\cel{3}gipso, cemento. - II. Petro-bloko kavigita o ligno-}
\l{\cel{3}vazo de qua drinkas la kavali, la bruti. - DE.}
\l{}
\l{\sou{troglodito}. I. Habitanto di la kaverni subsula. - II.}
\l{\cel{3}(\sou{zool}.)Ucelo insekto-manjanta, quan onu konfundas}
\l{\cel{3}multa-kaze a la regolo. -\cel{2}L.\cel{1}\sou{troglodytes}.- DEFIS.}
\l{}
\l{\sou{trokantero}. (\sou{anat.}) Singla ek la du tuberi rondatra di}
\l{\cel{3}la femuro ube (an la hancho) interjuntas la kolo}
\l{\cel{3}di la femuro e la korpo. - EF.}
\l{}
\l{\sou{trokaro}. (\sou{kirurg.}) Instrumento di qua la formo esas olta}
\l{\cel{3}di puncilo cilindra muntita an mancho e kontenata}
\l{\cel{3}en kanulo, quan onu uzas por puncionar. - DEFIS.}
\l{}
\l{\sou{trokeo}. (\sou{versif}.) Pedo ek un silabo longa ed un silabo}
\l{\cel{3}kurta. - DEFIR.}
\l{}
\l{\sou{trokoido}. (\sou{matem.})\cel{2}La figuro kurva quan prizentas kordo}
\l{\cel{3}fibranta, fuzelo. - DEFIS.}
\l{}
\l{\sou{troleo}. (\sou{tekn.}) Che la vehili elektro-movata : organo}
\l{\cel{3}por la preno di elektro-fluo, qua konsistas ye stango}
\l{\cel{3}flexebla kun mikra roteto o kun kontaktilo glitanta,}
\l{\cel{3}qua uzasas por transmisar la elektro-fluo de la kablo}
\l{\cel{3}konduktiva til la motoro di la vehilo. - DEFS.}
\l{}
\l{\sou{trombo}. Kolono de aquo qua, levata da vento-vortico, jir-}
\l{\cel{3}adas, devastante omno ube lu pasas. - DFIS.}
}
